\section{Discussion}
\label{sec:discussion}

\subsection{Advantages and Applicability Boundaries of Hypergraph Modeling}

The experimental results from the flow battery case reveal that hypergraph modeling delivers the most pronounced gains on three categories of high-complexity queries: multi-constraint problems~(Level~3), cross-material comparative questions~(Level~4), and strategy recommendation tasks~(Level~6). For multi-constraint problems, higher-order hyperedges preserve operational conditions and performance indicators within a single relational structure, mitigating the fragmentation inherent in pairwise graphs where condition--indicator pairs are scattered across separate edges. For cross-material comparisons, hypergraphs connect relation clusters belonging to different material systems through shared entities, reducing the reliance on manual alignment for lateral comparison. For recommendation tasks, higher-order relations sustain the semantic integrity of the ``objective--constraint--solution--expectation'' chain, yielding more coherent variable combinations in the generated recommendations.

These advantages are, however, bounded. When a query merely asks for the definition of a single concept, a simple factual lookup, or an enumerative overview of the corpus, vector retrieval or pairwise graph retrieval already achieves satisfactory accuracy. In such scenarios, the marginal gain from hypergraph modeling is marginal at best. The decisive factor is whether the query requires a judgment about ``whether multiple conditions belong to the same experimental fact''---only then does the structural information encoded in higher-order hyperedges become indispensable.

The comparison between the two cases reinforces this pattern. In the flow battery domain, higher-order relations primarily encode experimental-operation combinations linking ``system--active species--membrane/electrode--operating conditions--performance metrics.'' In the PFAS degradation domain, the corresponding pattern is ``target pollutant--catalytic material--reaction conditions--performance metrics--mechanism evidence.'' Notably, in PFAS Level~1 overview questions~(e.g., Q1 and Q2), the quality gap between the two systems is small because such questions depend mainly on enumerative knowledge with limited demands for condition completeness. By contrast, in Level~3 condition-combination questions~(Q5) and Level~5 mechanism-evidence-chain questions~(Q10), Hyper-ChE shows the largest advantage, precisely because these queries require assessing whether multiple variables co-occur within the same experimental fact or whether heterogeneous evidence jointly supports a single mechanistic claim. This cross-case consistency clarifies the applicability boundary: hypergraph modeling excels when multi-variable co-occurrence within a unified relational structure is essential to the query.

\subsection{Necessity of Domain Adaptation}

Generic LLMs face three distinct challenges when extracting entities from chemical-engineering literature. First, professional terminology and synonymy: the same material or concept may appear under multiple formulations~(e.g., ``Nafion membrane'' vs. ``perfluorosulfonic acid membrane,'' ``PFOA'' vs. ``C8''). Without domain constraints, an LLM may split synonymous references into distinct nodes. Second, non-standardized numerical conditions: temperature, concentration, and current density are reported in varying units or ranges, requiring structured parsing for normalization. Third, ionic variants: V(II)/V(III)/V(IV)/V(V) represent chemically distinct oxidation states in flow batteries, while PFAS species exhibit chain-length and functional-group variations that generic models struggle to classify correctly.

We address these challenges through two complementary measures. The first layer constrains the extraction space by pre-defining entity types and hyperedge schemas, guiding the LLM toward structured outputs over a limited type space and thereby reducing terminological ambiguity. The second layer enforces a JSON-structured output format with explicit field names and hierarchical relations, mitigating format drift.

It is worth emphasizing that no fine-tuning of the LLM is performed; all domain adaptation is achieved through prompt engineering. This design choice preserves cross-domain transferability: switching from flow batteries to PFAS degradation requires only replacing the entity-type and relation-type definitions, while the underlying extraction and retrieval pipeline remains unchanged. Based on the current corpus, this lightweight adaptation strategy has demonstrated basic viability across both chemical-engineering sub-domains.

\subsection{Limitations and Future Work}

The framework has two principal limitations that should be acknowledged frankly.

First, entity-name normalization remains incomplete. Different papers may refer to the same material with alternative formulations~(e.g., ``Nafion'' vs. ``perfluorosulfonic acid resin,'' ``PFOA'' vs. ``C8,'' ``PFOS'' vs. ``C8S''). The current framework does not integrate an external ontology for semantic normalization; handling such variants relies primarily on the LLM's in-context understanding. Although the type annotations in the pre-defined schema provide partial guidance, fully eliminating terminological ambiguity will require the introduction of domain-specific ontologies.

Second, the quality of higher-order hyperedge extraction exhibits stability issues. The structural correctness of a hyperedge depends on the LLM's comprehension of complex semantic relationships. When literature descriptions involve cross-sentence dependencies, implicit conditions, or intricate logical relations, the extracted hyperedges may suffer from entity omission or relation-type misclassification. In the flow battery case, the UNKNOWN rate stands at approximately 10--15\%, indicating that a non-negligible fraction of terms or relations fall outside the coverage of the pre-defined schema. Future work could improve stability through fine-tuning strategies or multi-round consistency verification---for instance, employing a verifier LLM to check the structural plausibility of extracted hyperedges.

Three directions merit further investigation: (i)~expanding corpus coverage to validate scalability; (ii)~developing hyperedge weight-learning mechanisms to optimize retrieval ranking; and (iii)~extending the method to additional chemical-engineering sub-fields such as electrocatalysis, battery materials, and reaction engineering to test cross-domain generality.

\section{Conclusion}
\label{sec:conclusion}

Analysis of literature across two representative chemical-engineering domains---flow batteries and PFAS degradation---reveals that the ``system--material--condition--metric--mechanism'' structure of experimental chemical knowledge is intrinsically multi-variate and coupled, making it naturally amenable to hypergraph modeling. To test this proposition, we developed Hyper-ChE, a domain-adapted hypergraph-enhanced RAG framework that replaces traditional pairwise graph structures with hypergraphs for knowledge storage and retrieval. Comparative experiments against Graph-RAG confirm that hypergraph modeling preserves condition-combination integrity more effectively than pairwise representations.

In the flow battery domain, a knowledge base constructed from 50 papers comprises 9,970 entities and 12,987 hyperedges, of which 20.9\% are higher-order. Hyper-ChE retrieves an average of 103 higher-order hyperedges and 63 entities per query, compared to 46 pairwise edges for Graph-RAG. In the PFAS domain, the corresponding knowledge base contains 7,476 entities and 10,812 relations/hyperedges with 21.98\% higher-order hyperedges. Across 12 test questions spanning overview, condition combination, cross-strategy comparison, long- and short-chain contrast, mechanism evidence chain, and experimental-scheme recommendation, Hyper-ChE averages 81.5 higher-order hyperedges against Graph-RAG's 37.9 pairwise edges. Qualitative analysis shows that Hyper-ChE performs particularly well on multi-condition preservation, cross-strategy lateral comparison, target-pollutant differentiation, multi-evidence mechanism chaining, and structured experimental recommendations.

Taken together, the cross-case validation yields three concrete contributions. \emph{Condition completeness}: higher-order hyperedges retain the association among multiple conditions within a single experiment, reducing information loss caused by pairwise relation decomposition. \emph{Contextual interpretability}: structured organization of mechanism pathways and comparative relations furnishes retrieval results with both factual data and mechanistic explanation. \emph{Cross-domain associability}: a unified knowledge-modeling framework supports knowledge transfer across distinct chemical-engineering sub-fields. The domain differences between the two cases---electrochemical performance optimization in flow batteries versus pollutant-removal mechanisms in PFAS degradation---demonstrate that the method is not confined to a specific chemical-engineering direction and can extend to other multi-variate experimental sciences.

We acknowledge that each domain-specific hypergraph database currently embeds only 50 papers; while the results on this subset already indicate the advantages of hypergraph modeling for multi-variable combination queries, more comprehensive conclusions await validation on larger corpora. Future work will focus on scaling corpus coverage, learning hyperedge weights for retrieval ranking, and extending the framework to additional sub-domains.

\section*{Code Availability}
\label{sec:code}

The source code, pre-defined schemas, and constructed knowledge-base snapshots for both case studies are available at \url{https://github.com/[anonymized/Hyper-ChE]}.
